\documentclass[11pt]{article}
\usepackage[margin=0.88in]{geometry}
\usepackage[T1]{fontenc}
\usepackage{lmodern}
\usepackage{amsmath,amssymb}
\usepackage{xcolor}
\usepackage{hyperref}
\hypersetup{colorlinks=true,urlcolor=blue!55!black,citecolor=blue!55!black}
\newcommand{\HC}{\operatorname{HC}}
\newcommand{\Hol}{\mathcal H(\mathbb C)}
\title{Multiply generated hypercyclic algebras:\\
the question in arXiv:1710.10413 was answered by arXiv:1706.08651}
\author{Literature-status packet for run \texttt{fa\_banach\_001}}
\date{11 August 2026}

\begin{document}
\maketitle

\begin{center}
\fbox{\parbox{0.92\linewidth}{\textbf{Status: literature already answered
affirmatively.}  The source says that the existence of multiply generated
algebras of hypercyclic entire functions was unknown.  B\`es--Papathanasiou
give a countably generated \emph{free dense} hypercyclic algebra for the
differentiation operator, a much stronger existential answer.  This is a
literature identification, not a new construction from the present run.}}
\end{center}

\section*{Original question}

Luis Bernal-Gonz\'alez, J.~Alberto Conejero, George Costakis, and Juan
Seoane-Sep\'ulveda, \emph{Multiplicative structures of hypercyclic functions
for convolution operators}, arXiv:1710.10413v1 \cite{BCCS}, state on
source-PDF page~3:
\begin{quote}
So far, the existence of multiply generated algebras of hypercyclic entire
functions is unknown.
\end{quote}
Here a hypercyclic algebra for an operator $T$ on the Fr\'echet algebra
$\Hol$ is a subalgebra $A$ such that $A\setminus\{0\}\subseteq\HC(T)$.
The source then obtains an infinitely generated \emph{multiplicative group}
$G$ for subexponential convolution operators $\Phi(D)$, all of whose elements
except $1$ are hypercyclic.  It notes that the algebra generated by $G$ is
infinitely generated, but does not claim every nonzero sum in that algebra is
hypercyclic.  Thus the group theorem does not itself answer the algebra
question.

\section*{Independent decisive answer}

Juan B\`es and Dimitris Papathanasiou, \emph{Algebrable sets of hypercyclic
vectors for convolution operators}, arXiv:1706.08651 \cite{BP}, formulate on
source-PDF page~2 the precise question:
\begin{quote}
Is it possible to find a hypercyclic algebra for $D$ that is not finitely
generated?
\end{quote}
They immediately state that they answer it affirmatively.  Their
Theorem~1.1 proves, under a geometric level-set condition on a convolution
symbol $\Phi$, that a generic sequence in $\Hol^{\mathbb N}$ freely generates
a dense hypercyclic algebra for $\Phi(D)$.  Corollary~1.3 specializes this to
the MacLane differentiation operator $D$:
\[
  \HC(D)\cup\{0\}\quad\text{contains a dense free algebra on countably many
  generators.}
\]
This is strictly stronger than merely finding a multiply generated
hypercyclic algebra: the generators are algebraically independent, there are
countably many of them, and the resulting algebra is dense in $\Hol$.

\section*{Identification, chronology, and scope}

The match is exact at the existential level.  The source's unqualified
question asks whether any multiply generated hypercyclic algebra of entire
functions exists; the independent paper constructs one for $D$, itself a
convolution operator, and explicitly says it answers Aron's equivalent
not-finitely-generated formulation.  The supporting arXiv record was first
submitted in June 2017, whereas arXiv:1710.10413 was submitted in October
2017; its proof-bearing version was subsequently revised in March 2019.  The
chronology explains why this is best viewed as an independent literature
answer rather than a new result of this run.

The answer does not say that \emph{every} nonconstant subexponential symbol
$\Phi$ supports such an algebra.  That stronger universal question is not what
the quoted sentence asks, and the source itself only promises a multiplicative
group for arbitrary such $\Phi$.  Bounded searches of the exact quoted phrase,
``not finitely generated hypercyclic algebra,'' ``free hypercyclic algebra,''
the source title, and both arXiv identifiers found the B\`es--Papathanasiou
paper as the decisive direct match.  The packet stores both PDFs locally.

\begin{thebibliography}{9}
\bibitem{BCCS}
L. Bernal-Gonz\'alez, J. A. Conejero, G. Costakis, and J. Seoane-Sep\'ulveda,
\emph{Multiplicative structures of hypercyclic functions for convolution
operators}, J. Operator Theory 80 (2018), 213--228;
\href{https://arxiv.org/abs/1710.10413}{arXiv:1710.10413}.

\bibitem{BP}
J. B\`es and D. Papathanasiou,
\emph{Algebrable sets of hypercyclic vectors for convolution operators},
\href{https://arxiv.org/abs/1706.08651}{arXiv:1706.08651},
Theorem~1.1 and Corollary~1.3.
\end{thebibliography}

\end{document}
